% GENERATED FILE. Edit canonical lessons, then run npm run generate:books.
% canonical-source-hash: fnv1a64:70ec41fc6a2c53b0
% canonical-lessons: FR-C298-renverser, FR-C298-blesse, FR-C298-choc, FR-C298-airbag, FR-C298-cabosser

\chapter{Le choc}
\label{ch:fr-le-choc}
\hlchaptermodality{298}

\begin{chapteropening}
\textbf{By the end of this chapter:} \emph{I can describe a crash: knocking someone over, an injured person, the shock, the airbag and a dent.}

Complete the last lesson of chapter 298: I can describe a crash: knocking someone over, an injured person, the shock, the airbag and a dent.
\end{chapteropening}

\section[renverser]{renverser — to knock over, to run over; to spill}

\label{lesson:FR-C298-renverser}



\begin{quote}
Before the new one: say the French for the accident report, then the French for the skid.
\end{quote}

\subsection*{You'll want to know: renverser}
\textbf{renverser} — \textquotedblleft{}to knock over, to run over; to spill\textquotedblright{}.

\begin{quote}
Une voiture a renversé un enfant à vélo.
\end{quote}

\begin{quote}
\emph{A car knocked over a child on a bike.}
\end{quote}

Also to spill: \emph{renverser du café}.

\begin{cousinweb}
\emph{renverser} is built on Old French \emph{enverser} \textquotedblleft{}to turn upside down\textquotedblright{}, from \emph{envers} \textquotedblleft{}reverse\textquotedblright{}, which is inherited from Latin \emph{inversus}, from \emph{vertere} \textquotedblleft{}to turn\textquotedblright{} (Lewis \& Short).
\end{cousinweb}

\begin{grammarlens}[title={What you've built}]
One more word for telling, explaining and describing at B1.
\end{grammarlens}

\subsection*{Your turn}
\begin{itemize}
\raggedright
  \item \emph{Say it:} \emph{renverser}
  \item \emph{Say it:} \emph{renverser}, once more
  \item \emph{Say it:} a sentence of your own with \emph{renverser}
  \item \emph{Recall:} say the French for the accident report, then the French for the skid, then say \emph{renverser} again
\end{itemize}

\begin{tcolorbox}[breakable,colback=teal!4,colframe=teal!35!black,title={Before you move on}]
What does renverser mean? (\textquotedblleft{}To knock over, to run over; to spill\textquotedblright{}.) Say it once more.
\end{tcolorbox}
\section[le blessé]{le blessé — the injured person, the casualty}

\label{lesson:FR-C298-blesse}



\begin{quote}
Before the new one: say the French for the skid, then the French for to knock over.
\end{quote}

\subsection*{You'll want to know: le blessé}
\textbf{le blessé} — \textquotedblleft{}the injured person, the casualty\textquotedblright{}.

\begin{quote}
Les blessés ont été emmenés à l'hôpital.
\end{quote}

\begin{quote}
\emph{The injured were taken to hospital.}
\end{quote}

\emph{Un blessé grave}, \emph{un blessé léger} — a seriously, a slightly injured person.

\begin{cousinweb}
\emph{blessé} is the past participle of the verb meaning \textquotedblleft{}to wound\textquotedblright{}, which is usually traced to a Frankish word for \textquotedblleft{}to bruise\textquotedblright{}; the origin is not fully settled (EtymDB lists no parent).
\end{cousinweb}

\begin{grammarlens}[title={What you've built}]
One more word for telling, explaining and describing at B1.
\end{grammarlens}

\subsection*{Your turn}
\begin{itemize}
\raggedright
  \item \emph{Say it:} \emph{le blessé}
  \item \emph{Say it:} \emph{le blessé}, once more
  \item \emph{Say it:} a sentence of your own with \emph{le blessé}
  \item \emph{Recall:} say the French for the skid, then the French for to knock over, then say \emph{le blessé} again
\end{itemize}

\begin{tcolorbox}[breakable,colback=teal!4,colframe=teal!35!black,title={Before you move on}]
What does le blessé mean? (\textquotedblleft{}The injured person, the casualty\textquotedblright{}.) Say it once more.
\end{tcolorbox}
\section[le choc]{le choc — the shock, the crash}

\label{lesson:FR-C298-choc}



\begin{quote}
Before the new one: say the French for to knock over, then the French for the injured person.
\end{quote}

\subsection*{You'll want to know: le choc}
\textbf{le choc} — \textquotedblleft{}the shock, the crash\textquotedblright{}.

\begin{quote}
Le choc a été très violent.
\end{quote}

\begin{quote}
\emph{The crash was very violent.}
\end{quote}

\emph{Être en état de choc} — to be in shock.

\begin{cousinweb}
\emph{choc} is built on \emph{choquer} \textquotedblleft{}to strike against\textquotedblright{}, which is usually traced to Middle Dutch \emph{schokken} \textquotedblleft{}to jolt\textquotedblright{} (EtymDB lists no parent).
\end{cousinweb}

\begin{grammarlens}[title={What you've built}]
One more word for telling, explaining and describing at B1.
\end{grammarlens}

\subsection*{Your turn}
\begin{itemize}
\raggedright
  \item \emph{Say it:} \emph{le choc}
  \item \emph{Say it:} \emph{le choc}, once more
  \item \emph{Say it:} a sentence of your own with \emph{le choc}
  \item \emph{Recall:} say the French for to knock over, then the French for the injured person, then say \emph{le choc} again
\end{itemize}

\begin{tcolorbox}[breakable,colback=teal!4,colframe=teal!35!black,title={Before you move on}]
What does le choc mean? (\textquotedblleft{}The shock, the crash\textquotedblright{}.) Say it once more.
\end{tcolorbox}
\section[l'airbag]{l'airbag — the airbag}

\label{lesson:FR-C298-airbag}



\begin{quote}
Before the new one: say the French for the injured person, then the French for the shock.
\end{quote}

\subsection*{You'll want to know: l'airbag}
\textbf{l'airbag} — \textquotedblleft{}the airbag\textquotedblright{}.

\begin{quote}
L'airbag s'est ouvert pendant le choc.
\end{quote}

\begin{quote}
\emph{The airbag opened in the crash.}
\end{quote}

Official French says \emph{le coussin gonflable}, the usual word in Québec.

\begin{cousinweb}
\emph{airbag} was borrowed from English \emph{airbag} (EtymDB), a bag filled with air.
\end{cousinweb}

\begin{grammarlens}[title={What you've built}]
One more word for telling, explaining and describing at B1.
\end{grammarlens}

\subsection*{Your turn}
\begin{itemize}
\raggedright
  \item \emph{Say it:} \emph{l'airbag}
  \item \emph{Say it:} \emph{l'airbag}, once more
  \item \emph{Say it:} a sentence of your own with \emph{l'airbag}
  \item \emph{Recall:} say the French for the injured person, then the French for the shock, then say \emph{l'airbag} again
\end{itemize}

\begin{tcolorbox}[breakable,colback=teal!4,colframe=teal!35!black,title={Before you move on}]
What does l'airbag mean? (\textquotedblleft{}The airbag\textquotedblright{}.) Say it once more.
\end{tcolorbox}
\section[cabosser]{cabosser — to dent}

\label{lesson:FR-C298-cabosser}



\begin{quote}
Before the new one: say the French for the shock, then the French for the airbag.
\end{quote}

\subsection*{You'll want to know: cabosser}
\textbf{cabosser} — \textquotedblleft{}to dent\textquotedblright{}.

\begin{quote}
Le choc a cabossé la voiture.
\end{quote}

\begin{quote}
\emph{The crash dented the car.}
\end{quote}

\emph{Une voiture toute cabossée} — a car full of dents.

\begin{cousinweb}
\emph{cabosser} is usually linked to \emph{bosse} \textquotedblleft{}a bump\textquotedblright{}; the origin is not settled (EtymDB lists no parent).
\end{cousinweb}

\begin{grammarlens}[title={What you've built}]
One more word for telling, explaining and describing at B1.
\end{grammarlens}

\subsection*{Your turn}
\begin{itemize}
\raggedright
  \item \emph{Say it:} \emph{cabosser}
  \item \emph{Say it:} \emph{cabosser}, once more
  \item \emph{Say it:} a sentence of your own with \emph{cabosser}
  \item \emph{Recall:} say the French for the shock, then the French for the airbag, then say \emph{cabosser} again
\end{itemize}

\begin{tcolorbox}[breakable,colback=teal!4,colframe=teal!35!black,title={Before you move on}]
What does cabosser mean? (\textquotedblleft{}To dent\textquotedblright{}.) Say it once more.
\end{tcolorbox}
